% !TeX program = xelatex
\documentclass{resume}
\resumepagelimit{1}
\hypersetup{
  pdftitle={Graduate Admissions CV Template},
  pdfsubject={Research-focused master's and PhD applications}
}

% Start here: replace every bracketed prompt; delete entries you do not have.
% This is a one-page starting point, not a program's required format or length.
% Check the program's CV instructions. Change the page budget above if needed.
% Guide: https://github.com/deepdave98/latex-resume-and-cv-templates/blob/main/docs/graduate-admissions.md
% Use your actual role: a thesis or course project is not an RA appointment.
% List degrees newest first. Add a master's only if you have started/completed it.
% Optional grades: use the transcript's scale and the program's reporting rules.

\begin{document}

\resumename{Your Name}
\resumecontact{City, Region
  \contactsep\resumeemail{hello@example.com}
  \contactsep\resumelink{https://example.com}{Website}
  \contactsep\resumelink{https://github.com/your-handle}{GitHub}}
% Delete Website or GitHub if you have no relevant work to link.

\section{Education}

\jobentry{[Degree in Field]}{[University]}{[City, Country]}{Expected [Month Year]}
\begin{jobduties}
  \item Thesis: [Title]. Supervisor: [Name]. [In progress / completed Month Year].
  \item Selected coursework: [Two or three courses relevant to the program].
\end{jobduties}
% Remove the thesis line if you do not have one. Remove Expected after graduation.
% Duplicate this degree block for earlier degrees; do not duplicate thesis bullets.

\section{Research Experience}

\jobentry{[Thesis Student / Actual Research Role]}{[Department, University]}{}{[Month Year] -- Present}
\begin{jobduties}
  \item Investigate [specific question] using [method] and [data or materials].
  \item Developed [your part of the study]; compared [approaches or conditions] while controlling for [confounder].
  \item Observed [finding or limitation] in [scope tested]; checked [alternative explanation] using [analysis or experiment].
\end{jobduties}
% Use present tense only for ongoing work. If results are pending, state what
% you have completed and what remains untested instead of inventing a finding.

\jobentry{[Actual Research Role]}{[Lab, Institution]}{}{[Month Year] -- [Month Year]}
\begin{jobduties}
  \item Built [instrument, dataset, or implementation] for [research question], supervised by [Name].
  \item Validated [specific property] using [check]; documented [limitation] in [report, code, or lab handoff].
\end{jobduties}
% Delete this entry if your thesis is your only research experience.

% Optional publications: remove the leading "% " from the block below.
% Include only work you authored; preserve author order. Use your field's style.
% Keep submissions/drafts separate from accepted or published work. Do not list
% the same paper twice. Add a real DOI/URL only if you may share it.
% BEGIN OPTIONAL PUBLICATIONS
% \section{Publications}
% \begin{skillitems}
%   \item {[Authors in order]}. ``[Title].'' [Venue], [Year]. [Published / accepted].
% \end{skillitems}
% END OPTIONAL PUBLICATIONS

% A preprint is not an accepted paper. Replace the status accurately.
% BEGIN OPTIONAL PREPRINTS
% \section{Preprints}
% \begin{skillitems}
%   \item {[Authors in order]}. ``[Title].'' [Year]. Preprint, not peer reviewed.
% \end{skillitems}
% END OPTIONAL PREPRINTS

\section{Selected Academic Project}

\jobentry{[Course Project / Independent Study]}{[Project Title]}{}{[Month Year] -- [Month Year]}
\begin{jobduties}
  \item Tested [question] by implementing [your contribution] and comparing it with [reference method].
  \item Found [result or discrepancy] on [test conditions]; reported [what the comparison cannot establish].
\end{jobduties}
% Keep course work labeled. For a team project, describe your part, not the team's.

% Optional presentations: label the event and your contribution accurately.
% Attending a conference is not a presentation; a poster is not a paper.
% BEGIN OPTIONAL PRESENTATIONS
% \section{Presentations}
% \begin{skillitems}
%   \item ``[Title].'' Poster presented at [Event], [Month Year]. [Your contribution].
% \end{skillitems}
% END OPTIONAL PRESENTATIONS

% Keep this heading with its entry when optional sections push it onto page two.
\Needspace{8\baselineskip}
\section{Teaching and Service}

\jobentry{[Teaching / Mentoring / Service Role]}{[Department or Organization]}{}{[Month Year] -- [Month Year]}
\begin{jobduties}
  \item Led [tutorial or activity] on [topic]; prepared [materials] and helped [audience] with [specific difficulty].
\end{jobduties}
% Delete this section if it does not apply. Add awards only if received, naming
% the award, awarding body, and year. Do not invent selectivity or class rank.

\section{Research Skills}

\begin{skillitems}
  \item \textbf{Methods:} [Methods you have used and can explain].
  \item \textbf{Tools:} [Languages, software, or equipment used in your work].
\end{skillitems}

\end{document}
